% Compilen desde esta carpeta o desde la raíz del repositorio.
\makeatletter
\def\input@path{{../../preamble/}{book2026-2027/preamble/}}
\makeatother
\input{preamble.tex}

\usepackage[spanish]{babel}
\usepackage{lmodern}
\usepackage{titlesec}
\geometry{a4paper,margin=1.6cm}
\linespread{1.10}

\definecolor{HeadingColor}{HTML}{243B53}
\definecolor{InstructionColor}{HTML}{704214}
\definecolor{ConceptColor}{HTML}{205C50}
\definecolor{LinkColor}{HTML}{0047AB}
\hypersetup{allcolors=LinkColor}
\newcommand{\guidelink}[2]{\href{#1}{\textbf{\underline{#2}}}}
\newcommand{\instruction}[1]{\textcolor{InstructionColor}{\textbf{#1}}}
\newcommand{\concept}[1]{\textcolor{ConceptColor}{\textbf{#1}}}
\newcommand{\checkbox}{%
    \begingroup
    \setlength{\fboxrule}{0.4pt}%
    \setlength{\fboxsep}{2pt}%
    \fbox{\rule{0pt}{0.8ex}\rule{0.8ex}{0pt}}%
    \endgroup}

\setlength{\parskip}{4pt}
\setlength{\tabcolsep}{5pt}
\renewcommand{\arraystretch}{1.15}
\setlist{leftmargin=0.65cm,topsep=4pt,partopsep=0pt,parsep=0pt,itemsep=3pt}
\titleformat{\section}{\large\bfseries\color{HeadingColor}}{}{0pt}{}
\titlespacing*{\section}{0pt}{10pt}{3pt}
\newcolumntype{Y}{>{\raggedright\arraybackslash}X}

\hypersetup{pdftitle={C9--C10. Sinteticen el trabajo y preparen la feria},pdfsubject={Emprendimiento e Innovación, Grado 10, Periodo 2}}

\begin{document}

\begin{center}
{\Large\bfseries\color{HeadingColor} C9--C10. Sinteticen el trabajo y preparen la feria\par}
\vspace{3pt}
Emprendimiento e Innovación \enspace|\enspace Grado 10 \enspace|\enspace Periodo 2\\
Guía 5 de 5
\end{center}

\section*{Objetivo y punto de partida}
\concept{C9.} Coopera en el equipo aportando ideas, cumpliendo compromisos y resolviendo desacuerdos de manera constructiva.\\
\concept{C10.} Explica el valor diferencial de la solución frente a la necesidad inicialmente identificada.

Recuperen D1 y los H pertinentes (P1). Consulten la elección, el plan, las versiones, PR1 y los ajustes (P2, C2--C8). Revisen la bitácora desde la primera guía para reconstruir cómo colaboró el equipo.

\section*{Documento acumulativo e integridad de la evidencia}
Esta guía es de consulta. Continúen el \textbf{mismo documento del proyecto en LaTeX} en \guidelink{https://www.overleaf.com/}{Overleaf} o en su editor habitual. Actualicen la fecha y entreguen el PDF completo. Usen \textbf{P1} para el Periodo 1 y \textbf{P2} para el Periodo 2 al citar criterios. Enlacen los registros por código o referencia interna y conserven sus fuentes en la sección de origen.

Distingan \concept{evidencia} (lo observado o expresado), \concept{conclusión} (lo que infieren), \concept{supuesto} (lo pendiente de verificar) e \concept{idea} (una propuesta). Registren fuentes y fechas reales. Identifiquen toda simulación autorizada por el docente. \instruction{Si usan IA, limítenla al código LaTeX, al formato y a organizar contenido propio.} Evidencias, conclusiones, ideas, decisiones y aportes deben provenir de su trabajo.

\section*{C9 de P2. Sinteticen los aportes, compromisos y acuerdos}
\begin{enumerate}[series=pasos]
\item \textbf{Revisen las entradas de toda la secuencia.} Identifiquen aportes que influyeron en una alternativa, un prototipo, la prueba o un ajuste. Revisen qué compromisos se cumplieron, cuáles se reprogramaron y qué acuerdos permitieron avanzar.
\item \textbf{Construyan una síntesis verificable.} Para cada integrante, citen al menos un aporte y un compromiso con su resultado y código de bitácora. Expliquen cómo el equipo abordó un desacuerdo real: qué opciones consideró, cómo escuchó las razones y qué acordó. Si no registraron desacuerdos, indiquen cómo acordaron una decisión concreta.
\item \textbf{Vinculen la síntesis con el trabajo realizado.} Remitan a los registros donde se ve el efecto de un aporte o acuerdo, como una alternativa, una decisión de diseño o un ajuste. Registren en la bitácora los compromisos para el ensayo y quién responderá por cada uno.
\end{enumerate}

\concept{Evidencia de C9.} Una síntesis por integrante con aportes, compromisos y resultados trazables a la bitácora, más un acuerdo explicado y vinculado al proyecto.

\newpage
\section*{C10 de P2. Sustenten qué aporta la solución y ensayen}
\begin{enumerate}[resume=pasos]
\item \textbf{Elijan un referente concreto.} Comparen V3 con una alternativa ALT de P2, C1--C2, o con la forma actual de atender la necesidad, descrita en los registros de P1. Indiquen qué característica de V3 cambia la experiencia, para quién y cómo responde a D1.
\item \textbf{Expongan el respaldo y sus límites.} Citen un H para fundamentar la necesidad y una observación, un comentario de PR1 o un rasgo demostrable del prototipo para sustentar el aporte. Si la prueba se hizo con V2, aclaren qué se observó allí y qué esperan del ajuste en V3. Indiquen una limitación o una condición de viabilidad.
\item \textbf{Preparen y ensayen la presentación.} Consulten la \guidelink{file:checklist_evaluativo_feria.pdf}{lista de coevaluación de la feria}. Organicen el recorrido de abajo, elijan evidencias que puedan mostrar y distribuyan las intervenciones. Ensayen dentro del tiempo asignado, con participación equilibrada y respuestas sustentadas en sus registros.
\end{enumerate}

{\small
\begin{tabularx}{\linewidth}{|p{0.25\linewidth}|Y|}
\hline
\textbf{Momento} & \textbf{Qué pueden mostrar y explicar} \\
\hline
Problema y usuario & D1 y un hallazgo H que muestre a quién afecta la necesidad y cómo. \\
\hline
Solución y demostración & Acción del usuario con V3, un cambio entre versiones y la razón que lo motivó. Usen los apoyos visuales para aclarar ese recorrido. \\
\hline
Diferencia y alcance & Comparación concreta, respaldo, beneficiarios, condiciones de viabilidad y una limitación. Expliquen con sus palabras por qué eligieron la solución. \\
\hline
\end{tabularx}\par
}

Pídanle a otro equipo que escuche el ensayo y devuelva una fortaleza, una mejora concreta y una solución alternativa con una explicación de cómo funcionaría o por qué ayudaría. Anoten qué mejorarán antes de la feria.

\concept{Evidencia de C10.} Una explicación comparativa del valor diferencial con referentes internos, respaldo y límites, incorporada al guion de la presentación.

\section*{Resultado y continuidad}
Agreguen la síntesis de colaboración, la explicación del valor diferencial y el guion breve con los turnos. Usen ese guion y las referencias para mostrar el prototipo y responder preguntas durante la feria.

\section*{Verificación final}
\begin{itemize}[label=\checkbox]
\item La síntesis identifica aportes, compromisos y acuerdos con entradas de la bitácora.
\item Explicamos el valor frente a un referente concreto y distinguimos lo observado de lo esperado.
\item Ensayamos con los apoyos, el tiempo asignado y la participación de todo el equipo.
\item Compilamos y revisamos el PDF acumulativo con sus referencias.
\end{itemize}

\end{document}
